\documentclass[UTF8,a4paper,10.5pt]{ctexart}
\usepackage[left=1.65cm,right=1.65cm,top=1.35cm,bottom=1.35cm]{geometry}
\usepackage{enumitem,xcolor,hyperref,tabularx,array,titlesec}
\pagestyle{empty}
\setlength{\parindent}{0pt}
\setlist[itemize]{leftmargin=1.4em,itemsep=1pt,topsep=2pt,parsep=0pt}
\titleformat{\section}{\large\bfseries}{}{0pt}{}[\titlerule]
\titlespacing*{\section}{0pt}{7pt}{3pt}

\begin{document}
\begin{center}
{\LARGE\bfseries 姓名}\\[3pt]
联系方式 \quad | \quad 邮箱 \quad | \quad
\href{https://github.com/sylviaedward/SC-RAG}{GitHub: sylviaedward/SC-RAG}
\end{center}

\section{教育背景}
\textbf{南京理工大学} \hfill \textit{通信工程 本硕 \quad XXXX.XX--XXXX.XX}\\
美赛 M 奖；数学建模竞赛；ISCAS 论文投稿（按实际情况替换）

\section{求职意向}
上海/苏州 AI Agent 产品、应用开发或 RAG 工程实习

\section{项目经历}
\textbf{SC-RAG-Agent：面向随机计算文献的可追溯研究 Agent}
\hfill \textit{Python / FastAPI / Streamlit / Docker}\\
\textit{项目概述：}为随机计算研究者提供可追溯综述，避免证据丢失、比较偏置与引用幻觉。

\begin{itemize}
\item 纯向量检索难召回 P2LSG、Sobol 等缩写，改为 BM25+dense RRF 并加重排；覆盖 50 篇论文、固定 $k=60$，兼顾术语与语义。
\item 固定单轮检索常漏掉比较题目标，改为规则驱动补检；弃用 LLM planner 降低工具幻觉，按 alias 缺失/候选不足触发 2 类补检。
\item 模型会改写引用且证据不足仍硬答，改为 span-ID 契约、Pydantic 校验、四态门控和单次修复；5 类异常统一 fail-closed。
\item 只跑通难定位失败，改为 58题（50+8）基准、phase4/消融与 AgentTrace；FastAPI+Streamlit+Docker 交付 5 接口。
\end{itemize}

\section{技能}
\begin{tabularx}{\textwidth}{@{}>{\bfseries}p{2.4cm}X@{}}
技能 & Python / FastAPI / Streamlit / Docker / RAG（BM25+dense+RRF+rerank）/ 自研 agent loop / Prompt Engineering / Pydantic / pytest\\
\end{tabularx}

\section{补充信息}
根据岗位要求补充可实习时间、英语能力、城市与到岗安排。

\end{document}
